% Generated draft: requires mathematical and source review.
\documentclass[11pt]{article}
\usepackage[margin=25mm]{geometry}
\usepackage{fontspec}
\setmainfont{Latin Modern Roman}
\usepackage{amsmath,amssymb,mathtools,unicode-math}
\setmathfont{Latin Modern Math}
\usepackage{enumitem,needspace,xurl,hyperref}
\setlength{\parindent}{0pt}
\setlength{\parskip}{5pt}
\setlength{\emergencystretch}{3em}
\newcommand{\N}{\mathbb{N}}
\newcommand{\Z}{\mathbb{Z}}
\newcommand{\Q}{\mathbb{Q}}
\newcommand{\R}{\mathbb{R}}
\newcommand{\C}{\mathbb{C}}
\newcommand{\F}{\mathbb{F}}
\newcommand{\A}{\mathbb{A}}
\newcommand{\PP}{\mathbb{P}}
\newcommand{\E}{\mathbb{E}}
\newcommand{\eps}{\varepsilon}
\newcommand{\dd}{\,\mathrm{d}}
\begin{document}
\section*{UnsolvedMath ID: 30001800}
\textbf{Source problem code:} catalogue entry
\subsection{Definitions and mathematical statement}
Let \(p\) be a prime, let \(F\) be a field of characteristic \(p\), and let \(X\) be a smooth proper \(F\)-scheme, with structure morphism \(f\colon X\to\operatorname{Spec}F\). For an abelian group \(A\), write \(A_{\Q}=A\otimes_{\Z}\Q\). Let \(K_r(X)\) denote Quillen algebraic \(K\)-theory, with its graded multiplication, and let \(K_r^M(F)\) denote Milnor \(K\)-theory. Thus \(K_0^M(F)=\Z\), and the Milnor \(K\)-ring is generated by \(F^\times\) in degree one subject to the Steinberg relations. Write \(\sigma_r\colon K_r^M(F)\to K_r(F)\) for the symbol map.

The multiplicative-generation question in the supplied record asks whether, for every such \(F\) and \(X\), and every integer \(r\geq 0\), the map
\[
\mu_r\colon K_0(X)_{\Q}\otimes_{\Q}K_r^M(F)_{\Q}\longrightarrow K_r(X)_{\Q},
\qquad a\otimes b\longmapsto a\cdot f^*\sigma_r(b),
\]
is surjective. Equivalently for this formulation, is the graded ring \(K_*(X)_{\Q}\) generated by its degree-zero part together with the image of \(K_*^M(F)_{\Q}\)? No finiteness assumption on the field \(F\) is supplied.

The record also asks for the kernel of a purported surjection
\[
CH^n(X)_{\Q}\otimes K_{i-n}^M(F)_{\Q}
\twoheadrightarrow H_M^{i+n}(X,\Q(n)),
\]
where \(CH^n(X)\) denotes the codimension-\(n\) Chow group and \(H_M^a(X,\Q(b))\) denotes motivic cohomology in degree \(a\) and weight \(b\). The record does not specify the ranges of \(i,n\), the tensor-product base, or the construction of this arrow. In the nonnegative range \(i\geq n\geq0\), the usual motivic cup product instead has degrees
\[
H_M^{2n}(X,\Q(n))\otimes_{\Q}
H_M^{i-n}(F,\Q(i-n))
\longrightarrow H_M^{i+n}(X,\Q(i)).
\]
Consequently, the displayed target weight \(n\) is incompatible with that usual product when \(i\ne n\). The intended map and indexing require clarification; neither its asserted surjectivity nor the claimed equivalence with the generation question is adopted here as an established fact. Determining a kernel also requires specifying more information than surjectivity alone.
\subsection{Short English statement}
The main question is whether every rational algebraic \(K\)-class on every smooth proper scheme over a field of positive characteristic is a sum of products of degree-zero classes with classes coming from the field's Milnor \(K\)-theory. The accompanying motivic kernel question has unresolved indexing and map-definition issues. The supplied record labels the problem open; this draft does not independently verify that status.
\subsection{Sources}
\begin{itemize}
\item Supplied local UnsolvedMath registry record: permanent ID \(30001800\), title ``Multiplicative Generation of Rational K-Theory,'' source problem code \texttt{catalogue entry}, in the Oberwolfach Reports Open Problems collection. No individual proposer is identified; the supplied proposed year is 2011.
\item Supplied source citation: \emph{New Directions in Algebraic \(K\)-Theory} (2011), Oberwolfach Reports, pp.\ 1469--1509, DOI \texttt{10.4171/owr/2011/27}. Source URL and supplied literature-source link labelled ``original conjecture'': \url{https://doi.org/10.4171/owr/2011/27}. No precise page or passage identifying this problem is supplied.
\item Literature assessment dated exactly \texttt{2026-08-21}: supplied dated curation, not independently verified. The supplied status is ``open,'' and the supplied assessment says that a proposed perfect-pairing formulation of the Beilinson and Parshin conjectures is not known in the stated generality. No such pairing is specified in the supplied mathematical statement, so the connection of this assessment to the generation question remains unresolved. The background explicitly requests independent verification of time-sensitive status and cited claims.
\item The registry supplies the label \texttt{corrected\_verified} and reports a manual semantic, typographic, and LaTeX review of a meaning-preserving cleanup. These are supplied editorial assertions, not independent verification of the mathematical formulation or its relationship to the cited report.
\end{itemize}
\end{document}
